\chapter{Umsetzung}

\section{Gesamtanwendung}

\todo{Daniel}

\subsection{webifier Tests}
In diesem Kapitel wird der allgemeine Aufbau, welcher für alle Tests von webifier gilt, erläutert.

Um die Tests vom Gesamtsystem abzukapseln wird auf Docker gesetzt. Hierbei wird für jeden Test ein eigenes Image geschrieben. Die Tests werden vom Tester dann gestartet. So wird jeder Test in einem eigenen Container ausgeführt. So ist sichergestellt, dass die Tests unabhängig von äußeren Faktoren sind und sich gegenseitig oder das Gesamtsystem nicht beeinflussen.

Die Technologien der einzelnen Tests sind abhängig vom jeweiligen Test und werden deshalb in den jeweiligen Kapiteln erläutert. Die Ergebnisübermittelung der Tests an den Tester wird mittels JSON-Strings realisiert. Wie in Beispiel (...) zu sehen besteht das JSON aus dem Testergebnis und einer ResultInfo. Die ResultInfo varriert von Test zu Test. Hier können für jeden Test weitergehende Informationen übermittelt werden. Für den Test auf Portscanning wird beispielsweise eine Liste von verdächtigen Portanfragen übermittelt.

\begin{scriptsize}
\lstset{
    style=eclipsejavascript,
    caption={Result JSON},
    label={lst:resultjson}
}
\begin{lstlisting}
  {
  	"result": "clean" | "suspicious" | "malicious" | "undefined",
  	"info": {
  		...
  	}
  }
\end{lstlisting}
\end{scriptsize}

\begin{itemize}
  \item Beschreiben der Startparameter URL und ID
\end{itemize}

\subsection{webifier Tester}

\todo{Samuel}

\subsection{webifier Plattform}

\todo{Samuel}

\subsection{webifier Mail}

\todo{Daniel}

\subsection{webifier Data}

\todo{Samuel}

\subsection{webifier Statistics}
Webifier Statistics wird in R implementiert. Hierzu werden Flexdashboards\footnote{Siehe http://rmarkdown.rstudio.com/flexdashboard/index.html} verwendet. Zur Generierung der Grafiken wurde auf verschiedene Librarys, wie beispielsweise Plot.ly, zurückgegriffen um den Entwicklungsaufwand für die Visualisierungen zu minimieren. Die Anordnung der Grafiken wird über ein bestimmtes Layout definiert. Jede Grafik wird prinzipiell in 3 Schritten erstellt:

\begin{enumerate}
  \item Daten aus der MongoDB laden
  \item Daten in die benötigte Form transformieren
  \item Entsprechende API ansteuern für Generierung der Grafik
\end{enumerate}

\begin{scriptsize}
\lstset{
    style=eclipsejavascript,
    caption={Beispiel R-Grafik},
    label={lst:rgrafik}
}
\begin{lstlisting}
  ### Durchschnittliche Analysezeit

  ```{r}
  result <- dbGetQueryForKeys(mg1, 'webifierTestResultData',"{}", "{durationInMillis:1}",skip=0,limit=Inf)
  mean.dur <- mean(result$durationInMillis)/1000
  mean.dur <- round(mean.dur)
  tp <- seconds_to_period(mean.dur)
  valueBox(paste(minute(tp),'min ',second(tp),'s',sep=""), icon="fa-hourglass-half",color="grey")
  ```
\end{lstlisting}
\end{scriptsize}

Im Codebeispiel \ref{lst:rgrafik} ist der Codeablauf für eine Valuebox zu sehen. Dieses Beispiel wurde ausgewählt um den Erstellungsprozess für die Grafiken zu erklären. Dies lässt sich auf alle anderen Grafiken übertragen.

Die Überschriften der Grafiken werden mit \textit{\#\#\#} markiert. Der R-Code befindet sich in Chunks, diese werden speziell markiert um dem Compiler kenntlich zu machen welches der R-Code ist.

Im Beispiel werden zunächst benötigten Daten aus der MongoDB geladen. Da hier eine Valuebox für die Anzeige der durchschnittlichen Analysezeit generiert wird werden nur die Analysezeiten(durationInMillis) benötigt. Diese werden anschließend gemittelt und von Millisekunden in Minuten/Sekunden transformiert. Zur Erstellung der Valuebox muss nun nurnoch der Text, die Farbe und ein passendes Icon ausgewählt werden. Die Generierung und Platzierung übernimmt Flexdashboard. Als Ausgabe wird eine HTML-Datei generiert, welche dann in den Webserver eingebunden wird um sie für die Nutzer zugänglich zu machen.

\begin{figure}[H]
  \centering
  \includegraphics[width=5cm]{images/stats/valuebox}
  \caption{Generierte Valuebox}
  \label{fig:valuebox}
\end{figure}

In Abbildung \ref{fig:valuebox} ist die fertig generierte Valuebox mit Überschrift, Text und Icon in passender Farbe dargestellt.

Für stets aktuelle Grafiken wird das R-Skript für die Statistiken mehrfach täglich neu gebaut um die aktuellen Daten mit einzubeziehen. Von einer \textit{On the fly}-Generierung der Grafiken wurde abgesehen, da dies für den Server zu rechenintensiv wäre.

\section{Tests}

\subsection{Virenscan der Webseite}

\todo{Samuel}

\begin{itemize}
  \item Httrack (Umsetzung)
  \item Download aller Dateien der Webseite
  \item Scannen der Heruntergeladenen Dateien
  \begin{itemize}
    \item Clamav (Umsetzung)
    \item AVG (Umsetzung)
    \item CAV (Umsetzung)
  \end{itemize}
\end{itemize}

\subsection{Vergleich in verschiedenen Browsern}

\todo{Daniel}

\subsection{Überprüfung der Port-Nutzung}
Bei diesem Test wird überprüft ob die Seite versucht einen Portscan auf dem Computer des Anwenders zu betreiben. Hierfür werden 3 Techniken eingesetzt. Die wichtigsten Aufgaben werden von PhantomJS und Bro erledigt.
Bro ist ein Netzwerkmonitoringtool und wird hier genutzt um den Traffic welcher zwischen Webseite und Client entsteht zu protokollieren und in einer Logdatei abzuspeichern. PhantomJS ist ein \textit{headless Browser}, welcher genutzt wird um die Webseite aufzurufen und dessen Javascript auszuführen. Das ganze funktioniert hier ohne grafische Oberfläche.

Der Ablauf des Tests sieht wie folgt aus: Zunächst wird Bro intialisiert und es werden Filter angelegt um lediglich die Ports, der eingehenden Anfragen, mitzuloggen und in der Logdatei abzuspeichern. Ist Bro vollständig initialisiert und einsatzbereit startet PhantomJS mit dem Aufrufen der Seite und Ausführen des JavaScript-Codes. Währenddessen speichert Bro alle Netzwerkaktivitäten. Sobald der Durchlauf von PhantomJS abgeschlossen ist wird mittels Python die Validierung des Ergebnisses gestartet. Hier werden die angefragten Ports aus der Logdatei geladen und klassifiziert. Die Ports 80 und 443 werden verworfen, da diese die HTTP und SSL Ports sind und somit als harmlos klassifiziert werden können. Die weiteren Ports werden in einer Liste an riskanten Ports gespeichert. Die Anzahl an Ports in dieser Liste bestimmt nun das Ergebnis des Testes. Wurden keine verdächtigen Portanfragen gefunden wird das Ergebnis \textit{unbedenklich} übermittelt. Bei 1 oder 2 Ports in der Liste gibt der Test \textit{verdächtig} als Ergebnis zurück. Sollte die Anzahl größer gleich 3 sein wird die Seite von diesem Test als \textit{bedrohlich} eingestuft. Zusätzlich zum Ergebnis wird die Liste der riskanten Ports in der Ergebnisinformation weitergeleitet.

\subsection{Überprüfung der IP-Nutzung}
Der Test auf verdächtige IP-Anfragen ist bis auf 2 Änderungen identisch zu vorherigem Test auf Portscanning. Deshalb werden in diesem Kapitel nur die Unterschiede beleuchtet.

Der erste Unterschied liegt in der Initialisierung von Bro. Hier werden Filter angewendet um die IPs, der ausgehenden Anfragen, zu loggen. Hier müssen die ausgehenden Anfragen betrachtet werden, da bei dieser Art von Angriff versucht wird mittels clientseitig ausgeführtem JavaScript das Netzwerk des Anwenders auszuspähen. Den Aufruf der Seite übernimmt auch hier PhantomJS. Bei der darauf folgenden Validierung werden die IPs auf bekannte Heimnetzadressbereiche wie beispielsweise 192.168.178.* oder 192.168.2.* gemappt. Auch hier werden verdächtige IPs in einer Liste gespeichert. Die Anzahl der Elemente in dieser Liste bestimmt das Ergebnis des Testes. Hierbei sind die Schwellwerte identisch mit denen des Portscanning-Tests, also bei 0 Abfragen wird \textit{sauber} zurückgegeben, bei 1-2 wird \textit{verdächtig} zurückgegeben und bei >3 wird die Seite als \textit{bedrohlich} eingestuft. Zusätzlich zum Ergebnis wird die Liste der riskanten IPs in der Ergebnisinformation weitergeleitet.

\subsection{Prüfung aller verlinkten Seiten}

\todo{Daniel}

\subsection{Google Safe Browsing}

\todo{Daniel}

\subsection{Überprüfung des SSL-Zertifikats}

\todo{Samuel}

\subsection{Erkennung von Phishing}

\todo{Samuel}

\subsection{Screenshot der Seite}
Der Screenshot der Seite erfolgt über eine von PhantomJS gelieferte Methode um den Seiteninhalt aufzunehmen und als Bilddatei abzuspeichern. PhantomJS wird hierbei genutzt da der Test in einem Docker ohne grafische Benutzeroberfläche läuft und deshalb ein headless Browser nötig ist um die Seite aufzurufen. Nachdem die Seite in einer Bilddatei abgespeichert ist, wird diese als base64-encoded String weitergegeben. Der Test liefert immer das Ergebnis \textit{sauber}, welches aber für den Tester irrelevant ist, da der Screenshot-Test keine Gewichtung im Tester hat. In der Ergebnisinformation wird der base64 encodierte String weitergegeben, welcher dann von der Plattform interpretiert und als Bild für den Nutzer dargestellt wird.
